\section{Introduction}
Diffusion models \citep{sohl2015deep,ho2020denoising,song2020score} have recently shown significant success in generating high-quality samples, achieving state-of-the-art performance in tasks like image, video, and audio generation. During \emph{training}, they learn a \emph{score function}--the gradient of the log density of the data distribution convolved with noise--by reconstructing data that has been corrupted with Gaussian noise, a process called \emph{denoising score matching} \citep{vincent2011connection}. During inference, the learned score is applied iteratively to denoise Gaussian noise to clean samples.

Despite these successes, the learned score presents an apparent paradox: if training were perfect, the model would learn the \emph{empirical score function}--the score with respect to the finite training data--and simply replicate the training data at inference, never generating novel samples. Recent work has sought to resolve this paradox by proposing various mechanisms by which diffusion models \emph{underfit} the empirical score, due to inductive bias or smoothness of neural networks~\citep{kamb2024analytic,niedoba2024towards,scarvelis2023closed,pidstrigach2022score,yoon2023diffusion,yi2023generalization}.

In this work, we argue that understanding diffusion models requires us to investigate not just \emph{how} they underfit the empirical score, but \emph{where}. Our starting point is the mismatch between the distribution of training samples versus the distribution of inputs encountered during inference trajectories. Consider \Cref{fig:extrapolation}: during training, for most timesteps, samples concentrate in thin shells around each data point (\coloredul{xblue}{blue shells}). Within these shells, the learned score simply points back to the shells' centers, i.e., the training images (\coloredul{xxgreen}{green arrows}). In contrast, we find that during inference, denoising trajectories quickly land in regions \emph{beyond} these shells (\coloredul{xred}{red points}), \emph{which were effectively never supervised with the empirical score at training time}.

This observation prompts us to distinguish two distinct regions in the data space: the \xblue{\textbf{supervision region}}, where the model is supervised to approximate the empirical score during training, and the \xred{\textbf{extrapolation region}}, where the model receives no supervision but is queried at inference. Our scaling experiments show that as models improve, their approximation of the empirical score \emph{inside} the supervision region improves, while it worsens \emph{outside} (underfitting). In other words, underfitting does not occur in the supervision region--a property we term \textbf{selective underfitting}, which, as we demonstrate, \textbf{is essential for understanding diffusion models.}

We leverage selective underfitting to design a suite of controlled experiments that shed new light on diffusion models' \textbf{(1) generalization}, showing that the \emph{size of the supervision region} strongly affects a model's ability to generalize--enlarging this region degrades generalization, an effect overlooked in prior studies--and \textbf{(2) generative performance}, where our novel scaling law between generative performance and \emph{supervision loss} enables directly quantifying the model's ability to \emph{transfer} the supervised score to the extrapolation region. This framework not only clarifies the benefits of recent advances such as REPA~\citep{yu2024representation} and diffusion transformers~\citep{peebles2023scalable,ma2024sit}, but also suggests a general principle for designing better training recipes. Taken together, our findings establish \emph{selective underfitting} as a central principle for understanding diffusion models.

\textbf{Contribution.}~\xsienna{\textbf{(1)}} We demonstrate \emph{selective underfitting} in diffusion models across two theoretically grounded regions of data space, with empirical validation that refines prior, oversimplified views. Leveraging selective underfitting: \xsienna{\textbf{(2)}} We advance a mechanism for how diffusion models generate samples beyond their training data that is agnostic to architectural inductive biases, unlike previous theories. \xsienna{\textbf{(3)}} We propose a novel scaling law to quantify the efficiency of training recipes for diffusion models, clarifying why recent seminal approaches work and providing a unified perspective for designing efficient recipes.